% ===========================================================================
% Category: Clustering
% Generated from ML_Lab_Manual_Completed.md - edit the Markdown and re-run
% scripts/build_latex_manual.py instead of editing this file by hand.
% ===========================================================================

\chapter{Experiment 1: K-means
Clustering}\label{experiment-1-k-means-clustering}

\textbf{Category:} Clustering

\section{1. Objective}\label{objective}

Partition unlabeled observations into K clusters by iteratively
assigning points to the nearest centroid and updating centroids; select
K using the elbow and silhouette diagnostics.

\section{2. Dataset Used}\label{dataset-used}

Mall Customers (\passthrough{\lstinline!data/mall\_customers.csv!}), 200
customers, two features: Annual Income (k\$) and Spending Score (1-100).
Features are standardized (\passthrough{\lstinline!StandardScaler!}).

\section{3. Required Libraries}\label{required-libraries}

numpy, pandas, matplotlib, scikit-learn

\section{4. Brief Theory}\label{brief-theory}

K-means minimizes the within-cluster sum of squares J = Σ\_k
Σ\_\{x∈C\_k\} \textbar\textbar x − μ\_k\textbar\textbar². It alternates
assignment and centroid update until convergence (a local optimum).
Euclidean distance is scale-sensitive, so features are standardized. K
is chosen externally, e.g.~by the elbow (inertia) and silhouette
(cohesion vs separation).

\section{5. Procedure}\label{procedure}

\begin{enumerate}
\def\labelenumi{\arabic{enumi}.}
\tightlist
\item
  Load the dataset and keep the two numerical features.
\item
  Standardize the features.
\item
  Try K = 2\ldots10; record inertia (WCSS) and silhouette for each K.
\item
  Choose K at the elbow / silhouette peak and fit the final model with
  \passthrough{\lstinline!init='k-means++'!},
  \passthrough{\lstinline!n\_init=20!},
  \passthrough{\lstinline!random\_state=42!}.
\item
  Visualize clusters and centroids.
\item
  Report silhouette, Davies-Bouldin and Calinski-Harabasz.
\end{enumerate}

\section{6. Reference Implementation}\label{reference-implementation}

\begin{lstlisting}[language=Python]
for k in range(2, 11):
    kmeans = KMeans(n_clusters=k, init='k-means++', n_init=10, random_state=42)
    kmeans.fit(X_mall_scaled)
    wcss.append(kmeans.inertia_)
    silhouette_scores.append(silhouette_score(X_mall_scaled, kmeans.labels_))
## 10. Complete Code Listing

*Full runnable program for this experiment, extracted from `notebooks/01_clustering_algorithms.ipynb` (executed; results above are its actual output).*

**Listing 1.1 — Core numerical and plotting libraries**

```python


# Experiment 2: Modified K-means

**Category:** Clustering

## 1. Objective
Implement the prescribed transparent modification of K-means: K-Means++ initialization plus an explicit outlier-distance check with a documented percentile threshold before final cluster interpretation.

## 2. Dataset Used
Mall Customers (200 × 2, standardized). The modification is applied on top of the same features.

## 3. Required Libraries
numpy, matplotlib, scikit-learn

## 4. Brief Theory
"Modified K-means" is not a single standard algorithm; the manual defines it precisely for reproducibility: initialize with K-Means++, fit K-means, compute each sample's distance to its assigned centroid, take a percentile threshold (97th percentile) and flag more distant samples as candidate outliers. It is an experimental variant that adds outlier awareness, not a replacement for robust clustering.

## 5. Procedure
1. Fit K-means with `init='k-means++'`, K = 5.
2. Compute d_i = ||x_i − μ_{c_i}|| for every sample.
3. Set threshold = 97th percentile of the training distances.
4. Flag samples with d_i > threshold as candidate outliers.
5. Plot centroids/clusters with flagged points highlighted.
6. Compare with the standard labels.

## 6. Reference Implementation
```python
km_mod = KMeans(n_clusters=5, init='k-means++', n_init=20, random_state=42)
mod_labels = km_mod.fit_predict(X_mall_scaled)
dist_to_centroid = np.linalg.norm(X_mall_scaled - km_mod.cluster_centers_[mod_labels], axis=1)
threshold = np.percentile(dist_to_centroid, 97)
outlier_flag = dist_to_centroid > threshold
\end{lstlisting}

\section{7. Results}\label{results}

{\def\LTcaptype{none} % do not increment counter
\begin{longtable}[]{@{}ll@{}}
\toprule\noalign{}
Item & Value \\
\midrule\noalign{}
\endhead
\bottomrule\noalign{}
\endlastfoot
Distance threshold (97th percentile) & 1.072 (standardized units) \\
Flagged candidate outliers & 6 of 200 (3\%) \\
Cluster labels & identical to standard K-means (same partition) \\
\end{longtable}
}

The flagged customers are the most weakly attached points of their
clusters --- typically boundary customers with unusual income/spending
combinations. The cluster structure itself is unchanged, which is
expected: the modification adds an outlier interpretation layer.

\section{8. Discussion and Conclusion}\label{discussion-and-conclusion}

The variant is useful when the analyst needs to know which observations
are poorly represented by any centroid (e.g., unusual customers).
Because the threshold is a documented percentile, the flag rate is
controlled and reproducible (3\% here). The limitation is that
distance-to-centroid is biased for elongated or density-varying
clusters, and the flagged points are not automatically errors.

\section{9. Answers to Viva Questions}\label{answers-to-viva-questions}

\begin{itemize}
\tightlist
\item
  \textbf{Why is it important to define a modification precisely?}
  Without a formal definition, ``modified K-means'' is ambiguous; a
  precise version (init method, threshold rule, percentile) is
  reproducible and testable.
\item
  \textbf{How does K-Means++ differ from random initialization?}
  K-Means++ selects initial centroids with probability proportional to
  squared distance to the nearest chosen centroid, spreading seeds out
  and reducing convergence to poor local optima.
\item
  \textbf{Why might distance-to-centroid fail for elongated clusters?}
  Centroid distance is isotropic; a point can be close to the centroid
  along the long axis yet far along the short axis, so elongated
  clusters can flag legitimate points and missing real outliers.
\end{itemize}

\begin{center}\rule{0.5\linewidth}{0.5pt}\end{center}

\section{10. Complete Code Listing}\label{complete-code-listing}

\emph{Full runnable program for this experiment, extracted from
\passthrough{\lstinline!notebooks/01\_clustering\_algorithms.ipynb!}
(executed; results above are its actual output).}

\textbf{Listing 2.1 --- Core numerical and plotting libraries}

\begin{lstlisting}[language=Python]
# ---- Core numerical and plotting libraries ----
import os
import numpy as np
import pandas as pd
import matplotlib.pyplot as plt
import seaborn as sns


# ---- Scikit-Learn clustering models ----
from sklearn.cluster import KMeans, AgglomerativeClustering
# ---- Feature scaling + unsupervised evaluation metrics ----
from sklearn.preprocessing import StandardScaler
from sklearn.metrics import silhouette_score, davies_bouldin_score, calinski_harabasz_score
# ---- SciPy hierarchical clustering utilities (linkage + dendrogram) ----
from scipy.cluster.hierarchy import dendrogram, linkage

# Styling setup (consistent look across all notebooks)
plt.style.use('seaborn-v0_8-whitegrid' if 'seaborn-v0_8-whitegrid' in plt.style.available else 'default')
plt.rcParams['font.sans-serif'] = 'DejaVu Sans'
plt.rcParams['figure.figsize'] = (10, 6)
plt.rcParams['figure.dpi'] = 110

# Fix the random seed so every run reproduces the same results
np.random.seed(42)
print("Libraries successfully imported!")
\end{lstlisting}

\textbf{Listing 2.2 --- Load the Mall Customers dataset}

\begin{lstlisting}[language=Python]
# ---- Load the Mall Customers dataset ----
df_mall = pd.read_csv('../data/mall_customers.csv')
print("Dataset Shape:", df_mall.shape)
print(df_mall.head(3))

# ---- Feature selection: two features so the clusters can be shown in 2D ----
X_mall = df_mall[['Annual Income (k$)', 'Spending Score (1-100)']].values

# ---- Standardize (mean = 0, std = 1): distance-based algorithms are scale-sensitive ----
scaler = StandardScaler()
X_mall_scaled = scaler.fit_transform(X_mall)
\end{lstlisting}

\textbf{Listing 2.3 --- Modified K-Means (manual variant): K-Means++ +
outlier-distance check}

\begin{lstlisting}[language=Python]
# ---- Modified K-Means (manual variant): K-Means++ + outlier-distance check ----
km_mod = KMeans(n_clusters=optimal_k, init='k-means++', n_init=20, random_state=42)
mod_labels = km_mod.fit_predict(X_mall_scaled)

# Distance of every sample to its assigned centroid
dist_to_centroid = np.linalg.norm(X_mall_scaled - km_mod.cluster_centers_[mod_labels], axis=1)
threshold = np.percentile(dist_to_centroid, 97)      # documented percentile
outlier_flag = dist_to_centroid > threshold

print(f"Distance threshold (97th percentile): {threshold:.3f}")
print(f"Flagged candidate outliers: {outlier_flag.sum()} of {len(X_mall)}")

# ---- Plot clusters with flagged candidate outliers highlighted ----
plt.figure(figsize=(9, 6))
for cluster_id in range(optimal_k):
    pts = X_mall[mod_labels == cluster_id]
    plt.scatter(pts[:, 0], pts[:, 1], s=45, edgecolors='k', alpha=0.8)
plt.scatter(X_mall[outlier_flag, 0], X_mall[outlier_flag, 1], s=180, facecolors='none',
            edgecolors='red', linewidths=2.5, label=f'Flagged outliers ({outlier_flag.sum()})')
plt.title("Modified K-Means: K-Means++ with Outlier-Distance Flagging", fontsize=13, fontweight='bold')
plt.xlabel("Annual Income (k$)", fontsize=12)
plt.ylabel("Spending Score (1-100)", fontsize=12)
plt.legend()
plt.tight_layout()
plt.show()
\end{lstlisting}

\chapter{Experiment 3: Hierarchical
Clustering}\label{experiment-3-hierarchical-clustering}

\textbf{Category:} Clustering

\section{1. Objective}\label{objective-1}

Perform agglomerative hierarchical clustering, interpret the dendrogram,
and compare Ward linkage with the other algorithms.

\section{2. Dataset Used}\label{dataset-used-1}

Mall Customers (200 × 2, standardized).

\section{3. Required Libraries}\label{required-libraries-1}

scipy, scikit-learn, matplotlib

\section{4. Brief Theory}\label{brief-theory-1}

Agglomerative clustering starts with one cluster per point and
repeatedly merges the closest pair. Ward linkage merges the pair that
increases total within-cluster variance the least: ΔW(A,B) = (n\_A
n\_B)/(n\_A+n\_B) \textbar\textbar μ\_A − μ\_B\textbar\textbar². The
dendrogram shows all merges; a horizontal cut determines the number of
clusters.

\section{5. Procedure}\label{procedure-1}

\begin{enumerate}
\def\labelenumi{\arabic{enumi}.}
\tightlist
\item
  Standardize the features.
\item
  Compute the Ward linkage matrix.
\item
  Plot the dendrogram (truncated for readability) and choose a cut
  level.
\item
  Fit
  \passthrough{\lstinline!AgglomerativeClustering(n\_clusters=5, linkage='ward')!}.
\item
  Evaluate with silhouette, Davies-Bouldin and Calinski-Harabasz.
\item
  Discuss linkage choice.
\end{enumerate}

\section{6. Reference Implementation}\label{reference-implementation-1}

\begin{lstlisting}[language=Python]
linkage_matrix = linkage(X_mall_scaled, method='ward')
dendrogram(linkage_matrix, truncate_mode='lastp', p=25, leaf_rotation=45, show_contracted=True)
agg_model = AgglomerativeClustering(n_clusters=5, metric='euclidean', linkage='ward')
agg_labels = agg_model.fit_predict(X_mall_scaled)
\end{lstlisting}

\section{7. Results}\label{results-1}

{\def\LTcaptype{none} % do not increment counter
\begin{longtable}[]{@{}ll@{}}
\toprule\noalign{}
Metric & Value \\
\midrule\noalign{}
\endhead
\bottomrule\noalign{}
\endlastfoot
Silhouette Score & 0.5538 \\
Davies-Bouldin Index & 0.5779 \\
Calinski-Harabasz Score & 244.41 \\
\end{longtable}
}

The dendrogram shows a clear large merge gap that supports cutting at
five clusters. The partition is nearly identical in quality to K-means
(silhouette difference \textless{} 0.001).

\section{8. Discussion and
Conclusion}\label{discussion-and-conclusion-1}

Ward linkage produces compact, balanced clusters and requires no
initialization, at the cost of O(n²) memory and the need to choose a
cut. Single linkage would chain across the space and complete linkage
would fragment; Ward is the appropriate default for compact clusters of
similar size.

\section{9. Answers to Viva
Questions}\label{answers-to-viva-questions-1}

\begin{itemize}
\tightlist
\item
  \textbf{What does a dendrogram show?} The full merge hierarchy: leaf
  order, merge heights (distances) and the structure of nested clusters;
  cutting at a height yields a flat partition.
\item
  \textbf{What is linkage?} The rule that defines the distance between
  two clusters (single = closest pair, complete = farthest pair, average
  = mean pair, Ward = variance increase).
\item
  \textbf{When is Ward linkage inappropriate?} With
  non-spherical/elongated clusters, unequal cluster sizes or outliers
  --- it assumes compact, similar-variance clusters.
\end{itemize}

\begin{center}\rule{0.5\linewidth}{0.5pt}\end{center}

\section{10. Complete Code Listing}\label{complete-code-listing-1}

\emph{Full runnable program for this experiment, extracted from
\passthrough{\lstinline!notebooks/01\_clustering\_algorithms.ipynb!}
(executed; results above are its actual output).}

\textbf{Listing 3.1 --- Core numerical and plotting libraries}

\begin{lstlisting}[language=Python]
# ---- Core numerical and plotting libraries ----
import os
import numpy as np
import pandas as pd
import matplotlib.pyplot as plt
import seaborn as sns


# ---- Scikit-Learn clustering models ----
from sklearn.cluster import KMeans, AgglomerativeClustering
# ---- Feature scaling + unsupervised evaluation metrics ----
from sklearn.preprocessing import StandardScaler
from sklearn.metrics import silhouette_score, davies_bouldin_score, calinski_harabasz_score
# ---- SciPy hierarchical clustering utilities (linkage + dendrogram) ----
from scipy.cluster.hierarchy import dendrogram, linkage

# Styling setup (consistent look across all notebooks)
plt.style.use('seaborn-v0_8-whitegrid' if 'seaborn-v0_8-whitegrid' in plt.style.available else 'default')
plt.rcParams['font.sans-serif'] = 'DejaVu Sans'
plt.rcParams['figure.figsize'] = (10, 6)
plt.rcParams['figure.dpi'] = 110

# Fix the random seed so every run reproduces the same results
np.random.seed(42)
print("Libraries successfully imported!")
\end{lstlisting}

\textbf{Listing 3.2 --- Load the Mall Customers dataset}

\begin{lstlisting}[language=Python]
# ---- Load the Mall Customers dataset ----
df_mall = pd.read_csv('../data/mall_customers.csv')
print("Dataset Shape:", df_mall.shape)
print(df_mall.head(3))

# ---- Feature selection: two features so the clusters can be shown in 2D ----
X_mall = df_mall[['Annual Income (k$)', 'Spending Score (1-100)']].values

# ---- Standardize (mean = 0, std = 1): distance-based algorithms are scale-sensitive ----
scaler = StandardScaler()
X_mall_scaled = scaler.fit_transform(X_mall)
\end{lstlisting}

\textbf{Listing 3.3 --- Build the linkage matrix and draw the
dendrogram}

\begin{lstlisting}[language=Python]
# ---- Build the linkage matrix and draw the dendrogram ----
# 'ward' linkage merges the pair of clusters that increases within-cluster variance the least.
linkage_matrix = linkage(X_mall_scaled, method='ward')

# The dendrogram shows every merge; we truncate it to the last 25 merges for readability.
plt.figure(figsize=(12, 5))
dendrogram(
    linkage_matrix,
    truncate_mode='lastp',   # show only the last p merged clusters
    p=25,
    leaf_rotation=45,
    leaf_font_size=10,
    show_contracted=True     # condensed leaves are shown as counts
)
plt.title("Hierarchical Clustering Dendrogram (Ward Linkage)", fontsize=14, fontweight='bold')
plt.xlabel("Cluster Sample Index / Merged Size", fontsize=12)
plt.ylabel("Euclidean Distance (Ward Threshold)", fontsize=12)
plt.axhline(y=10.0, color='r', linestyle='--', label='Cut-off Threshold (K=5)')  # horizontal cut => cluster count
plt.legend()
plt.tight_layout()
plt.show()
\end{lstlisting}

\textbf{Listing 3.4 --- Fit the final Agglomerative model (Ward linkage,
K = 5)}

\begin{lstlisting}[language=Python]
# ---- Fit the final Agglomerative model (Ward linkage, K = 5) ----
agg_model = AgglomerativeClustering(n_clusters=optimal_k, metric='euclidean', linkage='ward')
agg_labels = agg_model.fit_predict(X_mall_scaled)

print(f"Hierarchical Silhouette Score: {silhouette_score(X_mall_scaled, agg_labels):.4f}")
print(f"Hierarchical Davies-Bouldin Index: {davies_bouldin_score(X_mall_scaled, agg_labels):.4f}")
\end{lstlisting}

\chapter{Experiment 4: Fuzzy C-means}\label{experiment-4-fuzzy-c-means}

\textbf{Category:} Clustering

\section{1. Objective}\label{objective-2}

Cluster data with soft memberships so each observation can belong to
multiple clusters with different degrees; report the fuzzy partition
coefficient (FPC).

\section{2. Dataset Used}\label{dataset-used-2}

Mall Customers (200 × 2, standardized), K = 5, fuzzifier m = 2.0.

\section{3. Required Libraries}\label{required-libraries-2}

numpy, pandas, matplotlib (the c-means equations are implemented from
scratch; \passthrough{\lstinline!scikit-fuzzy!} is an equivalent library
implementation)

\section{4. Brief Theory}\label{brief-theory-2}

Fuzzy C-means minimizes J\_m = Σ\_i Σ\_k u\_ik\^{}m
\textbar\textbar x\_i − v\_k\textbar\textbar² with memberships u\_ik ∈
{[}0,1{]}, Σ\_k u\_ik = 1. It alternates: v\_k = Σ\_i u\_ik\^{}m x\_i /
Σ\_i u\_ik\^{}m and u\_ik = 1 / Σ\_j (d\_ik/d\_ij)\^{}(2/(m−1)). The
fuzzifier m controls softness (m → 1 becomes K-means). FPC = (1/N) Σ\_i
Σ\_k u\_ik² measures fuzziness (1 = hard, 1/K = maximally fuzzy).

\section{5. Procedure}\label{procedure-2}

\begin{enumerate}
\def\labelenumi{\arabic{enumi}.}
\tightlist
\item
  Standardize the data.
\item
  Initialize the membership matrix randomly (rows sum to 1).
\item
  Alternate centroid and membership updates until
  max\textbar ΔU\textbar{} \textless{} 1e-5.
\item
  Derive hard labels via argmax membership.
\item
  Report metrics and FPC; visualize memberships/hard partition.
\item
  Discuss ambiguous samples and m.
\end{enumerate}

\section{6. Reference Implementation}\label{reference-implementation-2}

\begin{lstlisting}[language=Python]
Um = U ** self.m
centers = np.dot(Um.T, X) / Um.sum(axis=0)[:, np.newaxis]
dist = np.linalg.norm(X[:, np.newaxis, :] - centers[np.newaxis, :, :], axis=2) + 1e-10
U = (1.0 / dist) ** (2.0 / (self.m - 1.0))
U = U / U.sum(axis=1, keepdims=True)
labels = np.argmax(U, axis=1)
\end{lstlisting}

\section{7. Results}\label{results-2}

{\def\LTcaptype{none} % do not increment counter
\begin{longtable}[]{@{}ll@{}}
\toprule\noalign{}
Metric & Value \\
\midrule\noalign{}
\endhead
\bottomrule\noalign{}
\endlastfoot
Silhouette Score (hard labels) & 0.5547 \\
Davies-Bouldin Index & 0.5722 \\
Fuzzy Partition Coefficient (FPC) & 0.6711 \\
Iterations to convergence & \textless{} 50 \\
\end{longtable}
}

FPC = 0.67 (between 1/K = 0.2 and 1.0) shows the partition is mostly
crisp but retains soft boundary memberships. Customers with membership
vectors close to (0.5, 0.5, \ldots) sit between clusters and are the
genuinely ambiguous ones.

\section{8. Discussion and
Conclusion}\label{discussion-and-conclusion-2}

FCM converges to the same hard partition as K-means here (silhouette
0.5547) because the Mall Customers clusters are well separated; the
added value is the membership matrix, which quantifies uncertainty at
boundaries. Larger m makes memberships softer (lower FPC) without
changing the hard structure much; smaller m (\textasciitilde1.1)
approaches K-means.

\section{9. Answers to Viva
Questions}\label{answers-to-viva-questions-2}

\begin{itemize}
\tightlist
\item
  \textbf{How is Fuzzy C-means different from K-means?} It assigns
  membership degrees instead of hard labels, so each point belongs to
  all clusters with weights that sum to 1; K-means is the limiting hard
  case.
\item
  \textbf{What does m control?} The fuzziness exponent: larger m spreads
  membership across clusters (softer), m approaching 1 gives crisp
  K-means-like assignments.
\item
  \textbf{What does a membership vector represent?} The degree to which
  one observation belongs to each cluster; near-uniform values indicate
  an ambiguous point on a cluster boundary.
\end{itemize}

\begin{center}\rule{0.5\linewidth}{0.5pt}\end{center}

\section{10. Complete Code Listing}\label{complete-code-listing-2}

\emph{Full runnable program for this experiment, extracted from
\passthrough{\lstinline!notebooks/01\_clustering\_algorithms.ipynb!}
(executed; results above are its actual output).}

\textbf{Listing 4.1 --- Core numerical and plotting libraries}

\begin{lstlisting}[language=Python]
# ---- Core numerical and plotting libraries ----
import os
import numpy as np
import pandas as pd
import matplotlib.pyplot as plt
import seaborn as sns


# ---- Scikit-Learn clustering models ----
from sklearn.cluster import KMeans, AgglomerativeClustering
# ---- Feature scaling + unsupervised evaluation metrics ----
from sklearn.preprocessing import StandardScaler
from sklearn.metrics import silhouette_score, davies_bouldin_score, calinski_harabasz_score
# ---- SciPy hierarchical clustering utilities (linkage + dendrogram) ----
from scipy.cluster.hierarchy import dendrogram, linkage

# Styling setup (consistent look across all notebooks)
plt.style.use('seaborn-v0_8-whitegrid' if 'seaborn-v0_8-whitegrid' in plt.style.available else 'default')
plt.rcParams['font.sans-serif'] = 'DejaVu Sans'
plt.rcParams['figure.figsize'] = (10, 6)
plt.rcParams['figure.dpi'] = 110

# Fix the random seed so every run reproduces the same results
np.random.seed(42)
print("Libraries successfully imported!")
\end{lstlisting}

\textbf{Listing 4.2 --- Load the Mall Customers dataset}

\begin{lstlisting}[language=Python]
# ---- Load the Mall Customers dataset ----
df_mall = pd.read_csv('../data/mall_customers.csv')
print("Dataset Shape:", df_mall.shape)
print(df_mall.head(3))

# ---- Feature selection: two features so the clusters can be shown in 2D ----
X_mall = df_mall[['Annual Income (k$)', 'Spending Score (1-100)']].values

# ---- Standardize (mean = 0, std = 1): distance-based algorithms are scale-sensitive ----
scaler = StandardScaler()
X_mall_scaled = scaler.fit_transform(X_mall)
\end{lstlisting}

\textbf{Listing 4.3 --- Fuzzy C-Means class (from scratch)}

\begin{lstlisting}[language=Python]
class FuzzyCMeans:
    """
    Fuzzy C-Means (FCM) clustering - implemented from scratch with NumPy.

    Each sample belongs to every cluster with a membership degree u_ik in [0, 1],
    and the memberships over the K clusters sum to 1 per sample.

    Parameters
    ----------
    n_clusters   : number of fuzzy clusters K
    m            : fuzzifier exponent (m > 1; larger m => fuzzier memberships)
    max_iter     : maximum number of alternating-optimization iterations
    tol          : convergence tolerance on the maximum membership change
    random_state : seed for reproducible random initialization
    """

    def __init__(self, n_clusters=5, m=2.0, max_iter=150, tol=1e-5, random_state=42):
        self.n_clusters = n_clusters
        self.m = m
        self.max_iter = max_iter
        self.tol = tol
        self.random_state = random_state

    def fit(self, X):
        rng = np.random.RandomState(self.random_state)
        n_samples = X.shape[0]

        # Step 1: initialize the membership matrix U randomly and normalize every row to sum to 1
        U = rng.rand(n_samples, self.n_clusters)
        U = U / U.sum(axis=1, keepdims=True)

        # Step 2: alternate between centroid and membership updates until convergence
        for iteration in range(self.max_iter):
            U_old = U.copy()

            # 2a: centroid update - membership-weighted mean of the data
            Um = U ** self.m                                        # u_ik^m
            centers = np.dot(Um.T, X) / Um.sum(axis=0)[:, np.newaxis]

            # 2b: compute the (n_samples x n_clusters) distance matrix to every centroid
            dist = np.linalg.norm(X[:, np.newaxis, :] - centers[np.newaxis, :, :], axis=2)
            dist = np.fmax(dist, 1e-10)  # avoid division by zero

            # 2c: membership update - closer points receive higher membership
            #     u_ik = 1 / sum_j (d_ik / d_ij)^(2/(m-1))
            inv_dist = 1.0 / dist
            power = 2.0 / (self.m - 1.0)
            inv_dist_p = inv_dist ** power
            U = inv_dist_p / inv_dist_p.sum(axis=1, keepdims=True)

            # 2d: stop when memberships stop changing (convergence)
            if np.max(np.abs(U - U_old)) < self.tol:
                break

        # Store results: final centroids, full membership matrix, and hard labels
        self.cluster_centers_ = centers
        self.u_ = U
        self.labels_ = np.argmax(U, axis=1)  # hard assignment = cluster with maximum membership
        return self


# Fit FCM with the same K as the other algorithms (m = 2 is the standard fuzzifier)
fcm = FuzzyCMeans(n_clusters=optimal_k, m=2.0, random_state=42)
fcm.fit(X_mall_scaled)
fcm_labels = fcm.labels_
fcm_centers = scaler.inverse_transform(fcm.cluster_centers_)

# ---- Evaluation of the hard partition induced by the fuzzy memberships ----
print(f"FCM Silhouette Score: {silhouette_score(X_mall_scaled, fcm_labels):.4f}")
print(f"FCM Davies-Bouldin Index: {davies_bouldin_score(X_mall_scaled, fcm_labels):.4f}")
fpc = np.mean(np.sum(fcm.u_ ** 2, axis=1))   # fuzzy partition coefficient (1/K fuzzy ... 1 hard)
print(f"FCM Fuzzy Partition Coefficient (FPC): {fpc:.4f}")
\end{lstlisting}
